% 附录 A —— 闭式解基准汇总

\chapter{闭式解基准汇总}
\label{app:benchmarks}

本附录把散落在各模块章的闭式解与已知解基准集中在一处，
供两类读者使用：需要快速判断"这套代码验证到什么程度"的人，
以及需要为新增模块挑选基准的人。

\section{按模块汇总}

{\small
\begin{longtable}{@{}p{0.075\textwidth} p{0.36\textwidth} p{0.17\textwidth} p{0.30\textwidth}@{}}
\toprule
\textbf{章} & \textbf{基准} & \textbf{判据} & \textbf{实测}\\
\midrule
\endfirsthead
\toprule
\textbf{章} & \textbf{基准} & \textbf{判据} & \textbf{实测}\\
\midrule
\endhead
\bottomrule
\endfoot
6 & 网格结构不变量（规模、覆盖、单调性、界面坐标、节点切片、接触次序） &
  逐条相等 & 全部通过\\
7 & D0 晶格常数对实验 & $\le1\%$ & GaN $0.345\%$/$0.328\%$，AlN $0.032\%$/$0.060\%$\\
7 & D1 弹性常数对实验 & $\le10\%$ & GaN $0.000\%$/$1.728\%$，AlN $2.778\%$/$1.877\%$\\
7 & D2 自发极化对实验 & $\le15\%$ & GaN $0.000\%$，AlN $0.000\%$\\
7 & D3 带隙（HSE 采纳值）对实验 & $\le5\%$ & GaN $0.877\%$，AlN $3.065\%$\\
7 & PBE 原始带隙不得被采纳 & 合并值 $\neq$ PBE 值 & 通过（原始值仍可查）\\
8 & Ambacher 极化面电荷（$x=0.25$） & $\in(1.10,1.30)\times10^{13}$\U{cm^{-2}} &
  $1.2020\times10^{13}$\U{cm^{-2}}\\
8 & 屏蔽严格按比例 & $10^{-6}$ 位 & 通过\\
8 & 压电极化为负；势垒侧 $p_{\text{sp}}$ 低于沟道侧 & 符号 & 通过\\
9 & 解析不动点的不变量（上界、$\Ef\ge E_0$、收敛） & 逐条 & 通过\\
9 & 面密度落在物理区间 & $4.0$--$12.0\times10^{12}$\U{cm^{-2}} & 通过\\
10 & 双介质平行板（无面电荷） & $<10^{-10}$\U{V} & $<10^{-10}$\U{V}（精确）\\
10 & 双介质平行板 $+$ 界面面电荷 & $<10^{-10}$\U{V} & $<10^{-10}$\U{V}（精确）\\
10 & $\Ec$ 算子 vs.\ $\varphi$ 形式独立离散化 & $<10^{-10}$\U{eV} &
  $2.9\times10^{-13}$\U{eV}\\
10 & 无限深方阱 $E_n=n^2\pi^2\hb^2/(2\meff L^2)$ & 二阶收敛 & 二阶收敛\\
10 & 谐振子 $E_n=\hb\omega(n+\tfrac12)$ & $<10^{-3}$\U{eV} & 通过\\
10 & 三角阱 Airy 零点 $E_n=\abs{\zeta_n}(\hb^2F^2/2\meff)^{1/3}$ & 二阶收敛 &
  $5.5\times10^{-5}\to3.4\times10^{-6}$\\
11 & 均匀电阻 $J=qN_D\mumu V/L$ & $<0.1\%$ & $\mathbf{<10^{-6}}$\\
11 & 均匀电阻两端电流一致 & 一致 & $<10^{-9}$\\
11 & 电流与偏压严格正比 & 比值 $=2$ & $2.0\pm10^{-6}$\\
11 & Bernoulli 函数 $B(x)=x/(\ee^x-1)$ & $10^{-12}$ 相对 & 通过（宽区间）\\
11 & 反射恒等式 $B(-x)=B(x)+x$ & 恒等 & 通过\\
11 & 通量四个偏导对中心差分 & 相对误差小 & 通过（含场依赖迁移率）\\
11 & 异质结台阶上的零通量 & $<10^{-12}$ & 通过\\
11 & 平衡态端点电流 & $<10^{-12}$\U{A/m} & $2.7\times10^{-15}$\U{A/m}\\
13 & 线性系统 & $\le2$ 步 & 通过（精确到 $10^{-10}$）\\
13 & 标量根 $x^2-2$ & 局部阶 $>1.7$ & 通过（收敛到 $\sqrt2$）\\
13 & 正确 Jacobian 被校验器确认 & $<10^{-8}$ & 通过\\
13 & 错一个元素被检出 & $>10^{-2}$ 且 \akey{ok=False} & 通过（\akey{worst_column} 正确）\\
14 & 块三对角解对稠密解（四组尺寸） & $<10^{-10}$ & 通过\\
14 & 单块退化为稠密解 & $<10^{-12}$ & 通过\\
14 & 正则化等价于平移对角块 & $<10^{-10}$ & 通过\\
\end{longtable}
}

\section{按"能证明什么"重排}

同一批基准若按 §\ref{sec:verify-criteria} 的五类判据重排，
可以看出验证覆盖的分布：

\begin{veriftable}{基准的类型分布}{tab:appA-kinds}
\vrow{闭式解（格式精度）}{检验离散格式}{
  第 10 章（4 项）、第 11 章（4 项）、第 14 章（3 项）}
\vrow{收敛阶（阶数正确）}{检验加密行为}{
  第 10 章（3 项）、第 11 章（1 项）}
\vrow{制造解（任意解上的精度）}{检验算子}{
  第 11 章（1 项，观测阶 $2.008/2.002/2.001$）}
\vrow{不变量与趋势}{检验必须成立的关系}{
  第 6、8、9、10、11 章}
\vrow{对实验（材料层）}{检验常数而非格式}{
  第 7 章（D0--D3 四级）}
\vnote{两类判据明显地不均衡：\textbf{格式侧的证据远多于器件侧}。
      这正是 \crefchap{ch:verification-methodology} 中"基准阶梯只走到 L1"
      的具体表现 —— 本表下半部分（对实测）几乎是空的。}
\end{veriftable}

\section{尚未使用的基准类型}

以下基准类型在方案中列出但在本文档当前版本中尚无用例，
记录在此以说明验证体系的\emph{空白}而非仅仅它的成果：

\begin{itemize}[leftmargin=1.7em,itemsep=0.28em,topsep=0.35em]
  \item \textbf{网格收敛指数（GCI）}：本项目的收敛阶断言用的是
        Richardson 三点公式，尚未按 GCI 规程给出离散误差的置信区间；
  \item \textbf{与独立实现的交叉对比}：除 \crefchap{ch:transport_dd}
        的 MMS 与 \crefchap{ch:sp1d} 的双路径 Jacobian 之外，
        尚无与第三方 TCAD 的对比（L3）；
  \item \textbf{不确定度量化}：无敏感性分析、无蒙特卡洛或多项式混沌。
\end{itemize}

\cleardoublepage
